% ==============================================================================
% Лабораторная работа #1: Реферат в LuaLaTeX
% Компилятор: LuaLaTeX
% ==============================================================================

\documentclass[14pt]{extreport}

% --- Поддержка языков и шрифтов (Задание 2) ---
\usepackage{fontspec}
\usepackage{polyglossia}
\setmainlanguage{russian}
\setotherlanguages{english}

% Основной шрифт документа
\setmainfont{Times New Roman}
\newfontfamily\cyrillicfont{Times New Roman}

% --- Настройка геометрии полей (Задание 6) ---
% Левое: 30 мм, Правое: 10 мм, Верхнее: 20 мм, Нижнее: 20 мм
\usepackage[left=30mm, right=10mm, top=20mm, bottom=20mm]{geometry}

% --- Базовое форматирование абзацев и интервалов (Задание 6) ---
\usepackage{indentfirst}               % Красная строка для каждого первого абзаца
\setlength{\parindent}{1.25cm}        % Размер абзацного отступа
\setlength{\parskip}{0pt}             % Нулевое расстояние между абзацами
\linespread{1.5}                      % Полуторный межстрочный интервал

% --- Нумерация страниц внизу по центру (Задание 6) ---
\usepackage{fancyhdr}
\pagestyle{fancy}
\fancyhf{}                            % Очистка стандартных колонтитулов
\cfoot{\thepage}                      % Номер страницы внизу по центру
\renewcommand{\headrulewidth}{0pt}    % Отключение верхней разделительной линии

% --- Пакеты для математики, таблиц, графики и ссылок ---
\usepackage{amsmath, amssymb, amsfonts}
\usepackage{graphicx}
\usepackage{tabularx}
\usepackage{hyperref}
\hypersetup{
    colorlinks=true,
    linkcolor=blue,
    urlcolor=blue,
    citecolor=blue
}

% ==============================================================================
% Метаданные титульной страницы (Задание 3)
% ==============================================================================
\title{Реферат на тему:\\[0.5cm] \textbf{Архитектура и эволюция систем компьютерной верстки: от TeX до LuaLaTeX}}
\author{Студент: Тоц Леонид Александрович \\ Группа: ИВТ-2}
\date{\today}

% ==============================================================================
% Тело документа
% ==============================================================================
\begin{document}

% Титульная страница (номер скрывается автоматически)
\maketitle
\thispagestyle{empty}

% Содержание
\tableofcontents
\newpage

% ------------------------------------------------------------------------------
\chapter{Введение}
% ------------------------------------------------------------------------------
Современная подготовка научных документов высокой типографической сложности неразрывно связана с развитием систем компьютерной верстки. Традиционные визуальные процессоры (WYSIWYG) нередко сталкиваются с ограничениями при работе со сложными математическими выкладками, перекрестными ссылками и масштабными библиографическими каталогами.

Целью настоящей работы является изучение теоретических основ и практическое освоение компилятора нового поколения \textbf{LuaLaTeX}, объединяющего мощь классического макропакета \TeX{} с нативной поддержкой шрифтов OpenType и встроенным интерпретатором языка программирования Lua.

% ------------------------------------------------------------------------------
\chapter{Эволюция и архитектура систем верстки}
% ------------------------------------------------------------------------------

\section{Исторические предпосылки создания TeX и LaTeX}
Система \TeX{} была создана выдающимся американским ученым Дональдом Кнутом в 1978 году. Основной задачей разработки было обеспечение бескомпромиссного качества верстки монографий и математических текстов. 

В 1980-х годах Лесли Лэмпорт создал надстройку над движком Кнута, получившую название \LaTeX{}. Главным концептуальным новшеством стало разделение смыслового содержания документа и правил его визуального оформления. Автор структурирует текст при помощи логических дескрипторов, а система берет на себя вычисление межсимвольных интервалов, переносов и расположения плавающих объектов.

\subsection{Специфика компилятора LuaLaTeX}
В отличие от классического движка \texttt{pdfTeX}, ориентированного на восьмибитные шрифтовые кодировки, движок LuaLaTeX предоставляет:
\begin{itemize}
    \item Полную нативную поддержку кодировки UTF-8;
    \item Возможность подключения системных шрифтов форматов TrueType и OpenType через пакет \texttt{fontspec};
    \item Продвинутую типографику для многоязычных документов с помощью пакета \texttt{polyglossia};
    \item Доступ к внутренним узлам синтаксического дерева документа через интерпретатор Lua.
\end{itemize}

\subsection{Сравнительный анализ систем}
Для наглядного сопоставления поколений компиляторов обратимся к сравнительной таблице~\ref{tab:comparison}.

\begin{table}[h!]
\centering
\begin{tabular}{|l|c|c|c|}
\hline
\textbf{Характеристика} & \textbf{pdfLaTeX} & \textbf{XeLaTeX} & \textbf{LuaLaTeX} \\
\hline
Базовая кодировка ввода & ASCII / T2A & UTF-8 & UTF-8 \\
\hline
Поддержка TrueType/OpenType & Ограниченная & Нативная & Нативная \\
\hline
Встроенный скриптинг & Нет & Нет & Язык Lua \\
\hline
Основной языковой пакет & \texttt{babel} & \texttt{polyglossia} & \texttt{polyglossia} \\
\hline
\end{tabular}
\caption{Сравнительная характеристика компиляторов семейства \TeX}
\label{tab:comparison}
\end{table}

% ------------------------------------------------------------------------------
\chapter{Элементы форматирования и математический аппарат}
% ------------------------------------------------------------------------------

\section{Использование математических окружений}
Набор сложных формул является сильной стороной \LaTeX{}. В качестве классического примера эквивалентности массы и энергии рассмотрим уравнение~\eqref{eq:einstein}:
\begin{equation}\label{eq:einstein}
E = mc^2
\end{equation}

Для многострочных математических выводов с выравниванием по операторам используется окружение \texttt{align}. Рассмотрим систему уравнений~\eqref{eq:system}:
\begin{align}\label{eq:system}
\nabla \cdot \mathbf{E} &= \frac{\rho}{\varepsilon_0}, \\
\nabla \cdot \mathbf{B} &= 0.
\end{align}

\section{Иллюстрации и внешние ссылки}
Официальный дистрибутив пакета \texttt{graphicx} поставляется с демонстрационными графическими заглушками (например, \texttt{example-image-a}), доступными во всех сборках TeX Live и Overleaf. Пример включения графики представлен на рисунке~\ref{fig:arch_demo}.

\begin{figure}[h!]
\centering
\includegraphics[width=0.45\textwidth]{example-image-a}
\caption{Пример интеграции графического объекта в окружение figure}
\label{fig:arch_demo}
\end{figure}

Более подробную документацию по движку и исходный код спецификаций можно найти на официальном веб-ресурсе проекта \href{https://www.luatex.org/}{LuaTeX.org}.

% ------------------------------------------------------------------------------
\chapter{Заключение}
% ------------------------------------------------------------------------------
В ходе выполнения лабораторной работы были освоены базовые принципы верстки научных документов в системе \LaTeX{} на базе современного движка LuaLaTeX. Был сформирован структурированный документ, содержащий титульный лист, автоматическое оглавление, перекрестные ссылки на таблицы, рисунки и математические формулы. Установленные параметры полей, полуторный межстрочный интервал и корректный абзацный отступ соответствуют требованиям оформления научно-исследовательских рефератов.

% ------------------------------------------------------------------------------
% Список литературы (Задание 5)
% ------------------------------------------------------------------------------
\newpage
\chapter*{Список литературы}
\addcontentsline{toc}{chapter}{Список литературы}

\begin{enumerate}
    \item Кнут Д. Э. Все про \TeX{} = The \TeX{}book. --- М.: ИД «Вильямс», 2003. --- 560 с.
    \item Львовский С. М. Набор и верстка в системе \LaTeX{}. --- 3-е изд., перераб. и доп. --- М.: МЦНМО, 2003. --- 448 с.
    \item Lamport L. \LaTeX{}: A Document Preparation System. User's Guide and Reference Manual. --- 2nd ed. --- Reading: Addison-Wesley, 1994. --- 272 p.
\end{enumerate}

\end{document}